\begin{proof}[Proof of \cref{obj:lem:point-version}]
\begin{enumerate}
\item For the selected arm regressions, write
\[
\mu_{a,P}(x)=E_P[Y\mid X=x,A=a],
\qquad a\in\{0,1\},\quad x\in[0,1].
\]
By \cref{obj:ass:contrast-lipschitz}, with \(L\) the common smoothness radius,
\[
|\tau_P(x)-\tau_P(x')|\le L|x-x'|,
\qquad x,x'\in[0,1].
\]
Thus \(\tau_P\) is Lipschitz and hence continuous. % lean: CausalSmith.Stat.PrivateCateRoughdesign.continuous_tau

\item For a Borel set \(B_{\mathrm{cov}}\subseteq[0,1]\), the density representation and \cref{obj:ass:density} give
\[
P_X(B_{\mathrm{cov}})
=\int_{B_{\mathrm{cov}}}f_P(x)\,dx,
\qquad
f_P(x)\ge\tfrac12
\quad\text{for Lebesgue-almost every }x.
\]
The density is therefore positive almost everywhere, so
\[
P_X(B_{\mathrm{cov}})=0
\quad\Longrightarrow\quad
\int_{B_{\mathrm{cov}}}dx=0.
\]
Consequently, equality \(P_X\)-almost everywhere implies equality Lebesgue-almost everywhere.

We record the continuous-function argument used here and below. Define the continuous projection
\[
c_{\mathrm{proj}}:\mathbb R\longrightarrow[0,1],
\qquad
c_{\mathrm{proj}}(u)=\max\{0,\min\{1,u\}\}.
\]
Let
\[
v_1,v_2:[0,1]\longrightarrow\mathbb R
\]
be continuous functions agreeing Lebesgue-almost everywhere. Their continuous extensions \(v_1\circ c_{\mathrm{proj}}\) and \(v_2\circ c_{\mathrm{proj}}\) agree almost everywhere on \([0,1]\), because
\[
c_{\mathrm{proj}}(x)=x,\qquad x\in[0,1].
\]
If their values differed at a point of \([0,1]\), continuity would give a relative neighborhood of that point on which they differed. Such a neighborhood has positive Lebesgue measure, including when the point is an endpoint. Hence
\[
v_1(c_{\mathrm{proj}}(x))
=v_2(c_{\mathrm{proj}}(x)),
\qquad x\in[0,1],
\]
and therefore \(v_1(x)=v_2(x)\) everywhere.

Now let \(v:[0,1]\to\mathbb R\) be continuous and satisfy \(v=\tau_P\) \(P_X\)-almost everywhere. The density argument transfers this equality to Lebesgue-almost everywhere, and the continuous-function argument yields
\[
v(x)=\tau_P(x),\qquad x\in[0,1].
\] % lean: CausalSmith.Stat.PrivateCateRoughdesign.continuous_tau_unique

\item We next establish the range of the selected arm means. By \cref{obj:ass:control-holder},
\[
|\mu_{0,P}(x)-\mu_{0,P}(x')|
\le L|x-x'|^{1/10},
\qquad x,x'\in[0,1],
\]
so \(\mu_{0,P}\) is continuous. The identity
\[
\mu_{1,P}(x)=\mu_{0,P}(x)+\tau_P(x),
\qquad x\in[0,1],
\]
then makes \(\mu_{1,P}\) continuous as well.

For \(a\in\{0,1\}\), define the arm propensity and the arm-specific covariate measure by
\[
e_{a,P}(x)=
\begin{cases}
1-e_P(x),&a=0,\\
e_P(x),&a=1,
\end{cases}
\qquad x\in[0,1],
\]
\[
\Lambda_{a,P}(B_{\mathrm{cov}})
=P\{X\in B_{\mathrm{cov}},A=a\},
\qquad B_{\mathrm{cov}}\subseteq[0,1]\ \text{Borel}.
\]
The selected propensity identity gives
\[
\int_{B_{\mathrm{cov}}}e_{1,P}(x)\,dP_X(x)
=\Lambda_{1,P}(B_{\mathrm{cov}}).
\]
For the control arm, partitioning \(\{X\in B_{\mathrm{cov}}\}\) by treatment gives
\[
P_X(B_{\mathrm{cov}})
=\Lambda_{0,P}(B_{\mathrm{cov}})
+\Lambda_{1,P}(B_{\mathrm{cov}}),
\]
and hence
\[
\int_{B_{\mathrm{cov}}}e_{0,P}(x)\,dP_X(x)
=P_X(B_{\mathrm{cov}})
-\int_{B_{\mathrm{cov}}}e_P(x)\,dP_X(x)
=\Lambda_{0,P}(B_{\mathrm{cov}}).
\]
By \cref{obj:ass:overlap}, both arm propensities satisfy
\[
\tfrac14\le e_{a,P}(x)\le1,
\qquad a\in\{0,1\},\quad x\in[0,1].
\]
Thus
\[
\tfrac14 P_X(B_{\mathrm{cov}})
\le
\int_{B_{\mathrm{cov}}}e_{a,P}(x)\,dP_X(x)
=\Lambda_{a,P}(B_{\mathrm{cov}}).
\]
In particular, a \(\Lambda_{a,P}\)-null Borel set is \(P_X\)-null and, by the preceding density argument, Lebesgue-null.

The selected regression identities and the binary range of \(Y\) give, for every Borel \(B_{\mathrm{cov}}\),
\[
0\le
\int_{B_{\mathrm{cov}}}\mu_{a,P}(x)\,d\Lambda_{a,P}(x)
=
\int_{\{X\in B_{\mathrm{cov}},\,A=a\}}Y\,dP
\le
\Lambda_{a,P}(B_{\mathrm{cov}}).
\]
The selected arm regressions are integrable under their respective arm measures. The set-integral order criterion therefore yields
\[
0\le\mu_{a,P}(x)\le1
\quad\text{for }\Lambda_{a,P}\text{-almost every }x.
\]
Indeed, the displayed inequalities applied to the measurable sets where the regression is negative or exceeds one exclude either violation on a set of positive arm measure. Transferring null sets as above gives the same bound Lebesgue-almost everywhere. Consequently,
\[
\mu_{a,P}(x)=c_{\mathrm{proj}}(\mu_{a,P}(x))
\quad\text{for Lebesgue-almost every }x.
\]
Both sides are continuous, so the continuous-function argument makes this equality pointwise. Since the projection takes values in \([0,1]\),
\[
0\le\mu_{a,P}(x)\le1,
\qquad a\in\{0,1\},\quad x\in[0,1].
\]
Applying these bounds at \(x_0=\tfrac12\) gives
\[
-1\le
\theta(P)=\mu_{1,P}(x_0)-\mu_{0,P}(x_0)
\le1.
\] % lean: CausalSmith.Stat.PrivateCateRoughdesign.theta_mem_Icc

\item We establish the conditional causal identity separately for each arm. All variables used below are measurable and bounded, so their conditional expectations are well defined. The arm-probability identities imply, for every Borel \(B_{\mathrm{cov}}\),
\[
\int_{\{X\in B_{\mathrm{cov}}\}}e_{a,P}(X)\,dP
=
\int_{B_{\mathrm{cov}}}e_{a,P}(x)\,dP_X(x)
=
\int_{\{X\in B_{\mathrm{cov}}\}}\mathbf1\{A=a\}\,dP.
\]
The characterization of conditional expectation by integrals over covariate events therefore gives
\[
E_P[\mathbf1\{A=a\}\mid X]=e_{a,P}(X)
\quad P\text{-almost surely}.
\]

Equality on all Borel covariate sets also identifies the arm measure as
\[
d\Lambda_{a,P}(x)=e_{a,P}(x)\,dP_X(x).
\]
Using this representation and the selected arm-regression identity,
\[
\begin{aligned}
\int_{\{X\in B_{\mathrm{cov}}\}}
e_{a,P}(X)\mu_{a,P}(X)\,dP
&=\int_{B_{\mathrm{cov}}}
e_{a,P}(x)\mu_{a,P}(x)\,dP_X(x)\\
&=\int_{B_{\mathrm{cov}}}\mu_{a,P}(x)\,d\Lambda_{a,P}(x)\\
&=\int_{\{X\in B_{\mathrm{cov}},\,A=a\}}Y\,dP\\
&=\int_{\{X\in B_{\mathrm{cov}}\}}\mathbf1\{A=a\}Y\,dP.
\end{aligned}
\]
Thus
\[
E_P[\mathbf1\{A=a\}Y\mid X]
=e_{a,P}(X)\mu_{a,P}(X)
\quad P\text{-almost surely}.
\]

By \cref{obj:ass:exchangeability}, projecting the potential-outcome pair onto either coordinate gives
\[
Y(a)\perp\!\!\!\perp A\mid X\quad[P],
\qquad a\in\{0,1\}.
\]
On \(\{A=a\}\), \cref{obj:ass:consistency} gives \(Y=Y(a)\) almost surely; on its complement the treatment indicator vanishes. Hence
\[
\mathbf1\{A=a\}Y
=\mathbf1\{A=a\}Y(a)
\quad P\text{-almost surely}.
\]
Conditional independence then factors the conditional expectation:
\[
\begin{aligned}
E_P[\mathbf1\{A=a\}Y\mid X]
&=E_P[\mathbf1\{A=a\}Y(a)\mid X]\\
&=E_P[\mathbf1\{A=a\}\mid X]\,
  E_P[Y(a)\mid X].
\end{aligned}
\]
Combining the preceding identities yields
\[
e_{a,P}(X)\mu_{a,P}(X)
=e_{a,P}(X)E_P[Y(a)\mid X]
\quad P\text{-almost surely}.
\]
Since \(e_{a,P}(X)\ge\tfrac14\), cancellation gives
\[
E_P[Y(a)\mid X]=\mu_{a,P}(X)
\quad P\text{-almost surely}.
\]
Applying this first to the treated arm and then to the control arm, and using conditional-expectation linearity for the bounded potential outcomes, gives
\[
\begin{aligned}
E_P[Y(1)-Y(0)\mid X]
&=E_P[Y(1)\mid X]-E_P[Y(0)\mid X]\\
&=\mu_{1,P}(X)-\mu_{0,P}(X)
=\tau_P(X)
\quad P\text{-almost surely}.
\end{aligned}
\] % lean: CausalSmith.Stat.PrivateCateRoughdesign.tau_conditional_causal_contrast

\item Finally, let \(P'\in\mathcal P\) satisfy \(P'^{\mathrm{obs}}=P^{\mathrm{obs}}\). Use the same arm-measure definition with \(P'\) in place of \(P\). Each arm measure is determined by the observed marginal, so
\[
\Lambda_{a,P'}(B_{\mathrm{cov}})
=P'^{\mathrm{obs}}\{X\in B_{\mathrm{cov}},A=a\}
=P^{\mathrm{obs}}\{X\in B_{\mathrm{cov}},A=a\}
=\Lambda_{a,P}(B_{\mathrm{cov}}).
\]
Thus \(\Lambda_{a,P'}=\Lambda_{a,P}\). The outcome integrals in the selected regression identities are likewise integrals of observed coordinates. For every Borel \(B_{\mathrm{cov}}\),
\[
\begin{aligned}
\int_{B_{\mathrm{cov}}}\mu_{a,P'}(x)\,d\Lambda_{a,P}(x)
&=\int_{\{X\in B_{\mathrm{cov}},\,A=a\}}Y\,dP'^{\mathrm{obs}}\\
&=\int_{\{X\in B_{\mathrm{cov}},\,A=a\}}Y\,dP^{\mathrm{obs}}\\
&=\int_{B_{\mathrm{cov}}}\mu_{a,P}(x)\,d\Lambda_{a,P}(x).
\end{aligned}
\]
Both regressions are integrable under this common arm measure. Uniqueness of integrable functions from their integrals over all measurable sets yields
\[
\mu_{a,P'}(x)=\mu_{a,P}(x)
\quad\text{for }\Lambda_{a,P}\text{-almost every }x.
\]
The null-set transfers established above make this equality Lebesgue-almost everywhere. Both regressions are continuous by their respective model conditions, so the continuous-function argument gives
\[
\mu_{a,P'}(x)=\mu_{a,P}(x),
\qquad a\in\{0,1\},\quad x\in[0,1].
\]
Evaluating the control and treated equalities at \(x_0\) and subtracting gives
\[
\theta(P')
=\mu_{1,P'}(x_0)-\mu_{0,P'}(x_0)
=\mu_{1,P}(x_0)-\mu_{0,P}(x_0)
=\theta(P).
\] % lean: CausalSmith.Stat.PrivateCateRoughdesign.theta_eq_of_observed_eq
\end{enumerate}
\end{proof}
